% Created 2024-01-02 Tue 09:42
% Intended LaTeX compiler: pdflatex
\documentclass[a4paper]{article}

\usepackage{../preamble}
\author{Nicky}
\date{\today}
\title{Simulating  Arrival Times for a Non-homogeneous Poisson Process}
\begin{document}

\maketitle
\tableofcontents

\section{Intro}
\label{sec:org8d9a843}

We provide an algorithm to generate arrival times of a non-homogeneous Poisson process (nhpp) when the rate function \(\lambda(s)\) is piecewise constant.
Specifically, we assume given two sequences \(\{\lambda_{j}\}_{j=0}^{\infty}\) and \(\{s_{j}\}_{j=0}^{\infty}\), with \(s_{0}=0\),  the \(s_{j}\) strictly increasing, and \(s_{k}\to \infty\) as \(k\to\infty\), and then we take \(\lambda(s) = \lambda_{j}\) for \(s\in [s_{j}, s_{j+1}]\).
An example would be a daily scheme of arrival rates of patients per hour at a hospital, where \(\lambda_{0} = 1\) from \([0, 8)\), \(\lambda_{1} = 10\) from \([8, 12)\), \(\lambda_{2} = 8\) from \([12, 18)\) and \(\lambda_{4}=3\) from \([18, 24)]\),  update \(\{s_{j}\}\) accordingly).
\section{The Maths}
\label{sec:orga91c10a}

One way to simulate a non-homogenous Poisson process with piece-wise constant rates is as follows. Consider 


The paper \href{https://web.ics.purdue.edu/\~pasupath/PAPERS/2011pasB.pdf}{Generating Nonhomogeneous Poisson Processes} by Raghu Pasupathy contains further background.
In particular, it explains that we need to find the smallest \(t\) such that \(\Lambda (s, t) = \int_{s}^{t} \lambda(y)\d y =x\), where \(x = -\log u\) and \(u\sim\Unif{[0.1]}\).
\subsection{A constant rate throughout}
\label{sec:org07316a0}

If the rate \(\lambda(x) \equiv \lambda\) is contant throughout, finding such \(t\) is simple because \(\Lambda(s, t) = \lambda (t-s)\) so that \(t= s + s/\lambda\).
This gives the following algorithm:
\begin{enumerate}
\item \(s = 0\).
\item \(\mineq \log(u)/\lambda, \quad u \sim \Unif{[0,1]}\).
\item Yield \(s\) and return to step 2.
\end{enumerate}
\subsection{Piecewise Constant}
\label{sec:orgedb2160}

Let us generalize this proceduce so that we can deal with  piecewise constant  \(\lambda(\cdot)\).
Noting that the sequence \(\{s_{j}\}\) places a grid on \(\R^{+}\), we define
\begin{align*}
j &= \max\{i: s_{i} \leq s\}, & k &=\min\left\{i : \Lambda(s, s_{i})  >x\right\},
\end{align*}
i.e., \(s_{j}\) is the point on the grid just below \(s\), and \(s_{k}\) is the smallest point on the grid such that the integral \(\Lambda(s, s_{k})\) exceeds \(x\).
For the algorithm, if \(s+x/\lambda_{j}< s_{j+1}\), then we can set \(s += x/\lambda_{j}\), and return \(s\) as the next arrival time.
However, if if \(s+x/\lambda_{j}\geq s_{j+1}\), we use the following idea:
\begin{equation*}
x = \Lambda(s, t) = \Lambda(s_{j}, s_{k}) - \Lambda(s_{j}, s) - \Lambda(t, s_{k}).
\end{equation*}
Since \(\lambda(\cdot)\) is \(\lambda_{j}\) on \([s_{j}, s]\) and \(\lambda(\cdot) = \lambda_{k-1}\) on \([t, s_{k})\),
\begin{align*}
\Lambda(s_{j}, s) &= \lambda_{j}(s-s_{k}),  & \Lambda(t, s_{k}) = \lambda_{k-1} (s_{k}-t).
\end{align*}
Thus, if we have \(s_{j}\) and \(s_{k}\), the next arrival time \(t\) after \(s\) is easily  solved from the equation
\begin{equation*}
\lambda_{k-1}(s_{k} -t ) = x +  \Lambda(s_{j}, s) - \Lambda(s_{j}, s_{k}).
\end{equation*}
That is,
\begin{equation}\label{eq:nh8}
t = s_{k } -  (\Lambda(s_{j}, s_{k}) - x -  \Lambda(s_{j}, s) )/\lambda_{k-1}.
\end{equation}
\section{The Algorithm}
\label{sec:org7cfc631}

With this idea, we obtain the following algorithm. Assume that \(\lambda_{0} > 0\) (Otherwise search for the first \(\lambda_{j}>0\), and start from there).

\begin{enumerate}
\item \(s=0\), \(j=0\)
\item \(s \mineq \log(u)/\lambda_{j}\)
\item If \(s < s_{j+1}\), yield \(s\) and return to step 2, else
\item \(T = \lambda_{j}(s_{j+1}-s)\).
\item While \(T<0\): \(j \pluseq 1\), \(T \pluseq \lambda_{j}(s_{j+1}-s_{j})\).
\item \(s = s_{j+1} - T/\lambda_{j}\).
\item Yield \(s\) and return to step 2.
\end{enumerate}
\section{The Python <Code}
\label{sec:orge7b1e9a}

\begin{minted}[]{python}
from typing import Iterator

import numpy as np


def nhpp_range(
    start: float,
    end: float,
    times: Iterator[float],
    rates: Iterator[float],
    gen: np.random.Generator = np.random.default_rng(3),
) -> Iterator[float]:
    rate = next(rates)
    down, up = next(times), next(times)
    while np.isclose(rate, 0):
        rate = next(rates)
        down, up = up, next(times)
    now = max(start, down)
    while True:
        now += gen.exponential() / rate
        T = rate * (up - now)
        while T < 0:
            rate = next(rates)
            down, up = up, next(times)
            T += rate * (up - down)
        now = up - T / rate
        if now > end:
            break
        yield now

\end{minted}
\section{The Tests}
\label{sec:orgf13b55a}

\begin{minted}[]{python}
from datetime import datetime, timedelta
from itertools import cycle, accumulate
from collections import Counter
import numpy as np
from nhpp import nhpp_range


def test_floats():
    gen = np.random.default_rng(3)

    start = 0
    end = 40
    periods = [1, 1, 3]
    rates = [0, 1, 10]

    times = []
    counter = Counter()
    counter[0] = 0
    for time in nhpp_range(
        start, end, accumulate(cycle(periods), initial=start), cycle(rates)
    ):
        times.append(time)
        counter[int(time) % 5] += 1
    print(counter)

    A = np.array(times)
    X = A[1:] - A[:-1]
    print(X.mean(), X.std())
    theoretical_mean_rate = sum(p * r for r, p in zip(rates, periods)) / sum(
        periods
    )
    print(1 / theoretical_mean_rate)


def test_datetimes():
    start = datetime(2023, 8, 9, 10, 0)
    end = datetime(2023, 8, 10, 18, 0)
    rates = list(range(13))
    periods = [2] * len(rates)

    time = start
    for delta in nhpp_range(
        0,
        (end - start).total_seconds() / 3600,
        accumulate(cycle(periods), initial=0),
        cycle(rates),
    ):
        time += timedelta(hours=delta)
        print(time)
    print(start, time)


if __name__ == "__main__":
    test_floats()
    test_datetimes()
\end{minted}
\end{document}
